\section{The Opposite of Identity}

\begin{quotex}
Abba Euloge one day was unable to hide his sadness.

— Why are you so sad, Abba? asked an old man.

— Because I am beginning to doubt the intelligence of the brothers concerning the great realities
of God. This is already the third time, having shown them a piece of linen on which I had drawn a small red dot and
asking them what they saw, they all replied, a ``small red dot" but never ``a piece of linen". 

\end{quotex}
\paragraph{Recapitulation}
Before getting to the main point, about which I've been pondering for some time, I want to draw a
few more red dots. There are two ways to read Gornahoor. The first is for those who know the principles of Tradition to
be true and the second is for those who are willing to consider what might logically and necessarily follow on the
assumption those principles are true. There are three main markers of Tradition:

\begin{enumerate}
\item Exoteric religion 
\item Esoteric or metaphysical dimension 
\item Social structure 
\end{enumerate}
We presume the existence of exoteric religions, but don't dwell on the multitudinous practices and
beliefs except perhaps for illustrative purposes. The second marker is the main focus since there is a common
metaphysics and understanding of symbols underlying traditional teachings. The third is the social structure that would
follow from the metaphysics: it would be hierarchical, ordered, and open to the transcendent dimension.

Yet, in some circles there is still persistent confusion, resulting from the emphasis on marker (1) and a misunderstood
(3). There is the belief that an exoteric religion can be chosen on a whim based on mere personal preference, such as
an attraction to certain myths, art forms, etc., even to the point of arbitrarily reclaiming the superficial elements
of a lost tradition as the ``real" thing.

The confusion about marker (3) also presumes that social structures can arbitrarily be imposed. With the opposition to
the so-called modern world evinced by Guenon, Evola, Coomaraswamy, etc., the presumption is that any opposition
whatsoever to certain aspects of the modern world is ipso facto ``traditional". Absolutely not, as only a certain
opposition has that character, viz., one that is hierarchical, ordered, and transcendent.

Ignorance of metaphysics is the fundamental cause of confusion, just as it is the fundamental point of a Guenon. The
``transcendental unity of religion" does not mean all religions are the ``same". Clearly, an atheist or a post-modern
thinker has absolutely no clue what ``transcendental unity" might mean, yet many self-described ``traditionalists" take
such comments seriously. To repeat, this is not the place to debate whether tradition is true, but rather to highlight
what it means. Anyone is free to deny it; i.e., to assume that reality is fundamentally chaos and that matter takes on
forms accidentally and at random. But it takes a special kind of fool to complain about this kaleidoscopic play of
forms, as he is bound to repeat them ad infinitum. Instead, we recommend for him an attitude ``amor fati".

\paragraph{Normal and Healthy Society}
For those of us in the West truly interested in Tradition, it behooves us to take a close look at what tradition has
looked like. All the early writers on this topic are agreed that for Western Europe, the Middle Ages represented
Tradition in each of the three markers. Hence, we should investigate what and how men thought at that time. Evola, in
his self-defense, claims that he believes what cultured men considered normal, healthy, and sane prior to the French
Revolution, at least in Latin Europe.

For example, here is an article by Plinio Corrêa de Oliveira\footnote{\url{http://www.tfp.org/tfp-home/plinio-correa-de-oliveira/philosophical-self-portrait-plinio-correa-de-oliveira.html}}, a man who has that same world view; this may help in
understanding it. He came to a similar understanding of historical developments as Guenon's,
although there is no direct link. This is an illustration of the ``transcendental unity." Now I understand that there
are the cultured despisers of the Western tradition who want no truck with such a viewpoint. They, therefore, should
substitute the word for their own tradition. Nevertheless, the essence will be the same; the new tradition will have an
identical metaphysical understanding; the symbols will be comparable.

To understand how things happen in the world, we need to know the three fundamental forces: Providence, Will, and
Destiny. These represent a deeper understanding of the three gunas of the Vedas. Destiny is the automatic process of
energy and matter that treats things indiscriminately and without differentiation. For example, the force of gravity
will be the same for man and for a rock. Providence is an upward pull that does not operate by brute force. This leaves
Will as the field open to Man. However, when men do not know their true will, the actions will be arbitrary and will
conflict with others.

Since the will follows the intellect, it knowing is necessary before willing and daring. Without knowledge of God, men
will not know how to follow Providence. Without knowledge of science, men will be at the mercy of the material forces.
And without knowledge of himself, his will is a tool of desire and ignorance.

So to understand the force of events, it is necessary to understand how men think and what their basic assumptions are.
Everything else follows. Of course, few men are consistent thinkers, so most men will deny the consequences of their
assumptions; that is why it is so difficult to point them out. Nevertheless, as time passes, consistency will indeed
win out. Most people know this intuitively, even if not explicitly; they seem to grasp the idea of the ``slippery
slope".

\paragraph{Pride and Sensuality}
As Mr. de Oliveira points out, there are two obstacles to finding the true will, namely pride and sensuality.It is
important to note that this has nothing at all to do with ``moralizing" or tied to a specific religious tradition, since
these obstacles can be precisely defined metaphysically and psychologically.

Ignorance may be of three kinds. There are the unintelligent, the bulk of men, who are simply incapable of metaphysical
principles. In a traditional society, they can rely on the social structure to teach them true things through stories,
symbols, etc.; they will benefit from the ``normal and healthy" beliefs of the leaders. Ignorance itself has a simple
cure: to acquire knowledge. Hence, those men capable of it should devote their leisure time to that acquisition.

Yet, there are the prideful, full of self-importance, who refuse to conform to true principles, preferring instead their
own opinions. Even if the spiritual and political authorities seem to be running on empty, and are not worth following,
the principles they should be expounding are still true. The problem is that if everyone is following his own opinion,
social chaos ensues, as is obvious today. The most organized opinions will be able to organize society along their own
lines.

If a man does not believe in a higher order, then it makes perfect sense that the purpose of the intelligence and the
will is to satisfy desire. Unfortunately, desire is the cause of suffering as desires can never be permanently
satisfied. (I am not speaking here of satisfying ``needs".) Sensuality will cloud a man's mind, and
we see this even among those considered to be ``traditionalists".



\flrightit{Posted on 2013-10-09 by Cologero }

\begin{center}* * *\end{center}

\begin{footnotesize}\begin{sffamily}



\texttt{Thorsten on 2013-10-11 at 11:19 said: }

Some important points have once again been need touched upon here that need resolution. It would be fascinating to find
out how the Traditionalist Christians differ from the self-styled ``Traditionalists" in the Alt Right and New Right in
terms of ethnic background. I'm confident that the latter group are overwhelmingly of Germanic (in
the broad sense of the term) and to a lesser degree Balto-Slavic descent, while the former have a much stronger Latin
component. Perhaps this is an all too obvious phenomenon but I think that in order for a Traditional society to take
shape again in the US or Europe the fact that the identitarian/antimodern Teutons can only be reeled in to Christianity
by making it resemble a pre-Christian warrior cult as much as creative license allows (Witness to what lengths the
authors of the Dream of the Rood and Heliand went in order to make the Gospel digestable.) This is something that will
have to be worked out if for no other reason than to get men of warrior spirit on the right side of this fight.

I myself belong to this nettlesome group and have been earnestly been trying to find a way to make traditional
Christianity work for me. But I cannot emphasize enough how truly foreign this religion feels, all the more so to the
degree that the Old Testament emphasized. Consider how much longer Christianity has settled into the makeup of the
culture of Latin Europe as opposed to Iceland or Lithuania for example where pre-Christian beliefs persisted into
modernity.

At the same time Traditionalists (of any faith) are by and large right in saying that Germanic, or for that matter any
other stripe of European, neo-paganism is contrived because it does indeed lack the degree of continuity in terms of
initiation and institutions. This does not negate the fact that many of these individuals truly descend from peoples in
which Christianity only ever achieved a relatively superficial penetration into the hearts of its coverts.

I think a dialogue is in order and I hope individuals from both camps can hear each other out for both sides have
ignored one another's disagreements and both sides will benefit from the exchange.


\hfill

\texttt{Cologero on 2013-10-12 at 10:51 said: }

Thorsten, you are bringing up the question of the proper vehicle for Tradition and how it may vary from people to
people. Although at the level of metaphysics, a ``dialog"should be possible, at the psychological level, the situation
is quite different. That is the domain of the psyche or soul, the vast middle ground between the body and spirit.
Unfortunately, Guenon simply ignores that, even though that is where people live out their lives. Here is where Jung is
of value, since his domain is the psychological effects of religious doctrines, not the truth of the doctrines
themselves.

It is obviously true, as you have noted, the Christianity is commonly taught today has lost its hold on most intelligent
men. In your case, as you say, it feels ``alien". Attempts to ``upgrade" it have had mixed results. On the one hand,
there are the demythologizers like a Marcus Borg (Nordic) who try to make it suitable for the modern mind. Although he
brings up many good points, his worldview is hardly distinguishable from the secular modern mind.

On the other hand, a Valentin Tomberg goes the other way to see the Christian message as a series of true myths, a
revelation of the hidden dimension of history. Presumably, he is not introducing novelty, but merely bringing attention
to what was already there. Unfortunately, his work is quite difficult, requiring years of study and meditation. It
remains to be seen if his work will be the leavening he hoped it to be.

That aside, Gornahoor is not trying to vivify Christianity, but rather to show that at one time its doctrines and
metaphysics were compatible with Tradition. That is no different from what a Guenon or a Coomaraswamy also claimed. In
the modern world, there are many choices, which are ancestors did not have, so there is the feeling (actually an
illusion) that we can simply choose a tradition much like a housewife picking out fruits at the marketplace.

The point I have been trying to make is that the overturning of Christian civilization has consequences, as it is
simultaneously tied to the rise of the modern world. However, one feels about Christianity personally, that is simply
an objective fact. So my gripe with that new right crowd is that they end up as unconscious enablers of the modern
world. Moreover, their criticisms are puerile and ignorant, usually derived from the perennial opponents of Christian
civilization.

The task, as I see it, for those who want both to cleave to tradition and to overcome Christian civilization is to be
constructive rather than destructive. Just as you pointed out, Christianity in Europe developed from, built on, and
integrated customs and insights from the previous pagan civilization; we have not denied that here. On the contrary, we
have pointed it out over and over. So the future neo-paganism, if there should ever be such a prophet who is yet to
appear, would have to do a similar thing.


\hfill

\texttt{Thorsten on 2013-10-12 at 23:41 said: }

Cologero,

I appreciate your response. I have not yet read Jung because I had gotten the impression that his work was incompatible
with Tradition. But given his value ``in the domain of the psyche or soul, the vast middle ground between the body and
spirit" I will have to put him on my list likewise with Tomberg.

Although I understood that Gornahoor has sought to, and to my understanding succeeded in demonstrating the compatibility
of medieval Christendom Tradition, I was admittedly confused about Gornahoor's position on
Christianity, having been under the impression that an attempt at the revival of Christianity in line with Tradition
was indeed a goal of Gornahoor.

So if Gornahoor is not attempting to vivify Christianity (and I would assume doesn't find the
vivification thereof possible or worthwhile), but if choosing one's tradition is not an option
then where does that leave those of the West who cannot find salvation within Christianity given that (neo-)paganism is
(at best) little more than various muddled attempts at reconstructing ancient Indo-European religious practices based
on a collection of mythological, philosophical, and ethnographic literature? Is the arrival of a prophet as you allude
to who could develop, build on, and integrate customs and insights from pagan antiquity and medieval Christendom (Btw I
didn't mean to suggest Gornahoor had in any way denied the influence of paganism on historical
Christendom) the only conceivable solution for those in the West who cannot honestly call Christendom home?

I don't mean to put you on the spot but I am struggling to find a workable solution for those who
are effectively unmoored from a living tradition. This possibility of a new prophet or a divine king is precisely what
I suggested was the only hope not just for the tradition-bereft but for the West as a whole in comment on a post on a
well known New Right website. To my disappointment, but hardly to my surprise this received no positive commentary.
However I think I succeeded in pointing out the spiritual hollowness of the worship of one's own
race. To believe one's race or nation has a divine mission or higher destiny can be positive in
certain contexts and has a traditional precedent at least as far back as Virgil, but equating
one's biological race with the divine is nothing more than collectivized narcissism (or pride and
sensuality) and needless to say is incompatible with Tradition. It really is a shame because there are people of great
intellectual and leadership ability, who offer valid critiques of modernity which relies upon an increasingly rare
affective detachment and insight, who seem to reject out of principle the possibility of metaphysical reality.

Regards,

Thorsten


\hfill

\texttt{Cologero on 2013-10-13 at 01:23 said: }

If you read Jung as a scientist and not as a metaphysician, I don't see a problem. I recommend you
read paragraphs 340 and 341 in the chapter Symbols of the Mother in Symbols of Transformation and then get back to me.

I don't know exactly what you are referring to, but it seems that most of that crowd has no
objection to the modern world per se, they want it populated just with white people. You can fight political battles on
certain issues and take allies of convenience, but that is not the real problem.

The modern world is based on the overturning of the once accepted order of things. That is beyond dispute and it
doesn't require any great intellectual ability to discern that. The moderns will freely admit it.
To oppose modernity therefore is to desire to restore the traditional order of things. That may sound simple enough,
but over the centuries most men have forgotten what that was, or else they mistake the faint echos of that order for
the real thing. I would recommend staying away from such men as their ``insights" are illusory, but you are free to
spend your time and efforts any way you like.

Now, Thorsten, you are putting yourself in quite a bind. You reject both the modern world and the spiritual source of
the order that the modern world overturned and continues to overturn. No wonder you feel frustrated and at a loss. 

No, Gornahoor is not ``trying" to revive Christianity along the lines of Tradition, just the opposite in fact. I agree
with Guenon and Coomaraswamy that medieval (in my definition of it) Christianity was traditional. All I have done is to
bring to light certain texts to demonstrate that. If you reject Tradition out of hand, as those ``great intellectuals"
do, then this whole project is meaningless and their criticisms of Christianity are really irrelevant. In effect, they
are also denying the source of the order that served their ancestors quite well. If they failed to absorb that teaching
fully — and I don't believe that — then they need to
yield to those who do understand it. Or do you believe that all opinions are of equal value?

If you claim to be a ``traditionalist", then you need to actually live that way. It is more than an intellectual
achievement as the moral component is even more important. It requires loyalty to one's ancestors
and to spiritual authority. It requires acceptance of a hierarchy. It requires adherence to the divine order of things.
It is one thing to say so, quite another to do so. I am not in a position to judge that, nor do I suffer from the
grandiose notion that the world is waiting for me to ``vivify" tradition. However, I do know that I would never under
any circumstances submit to those you claim have ``leadership ability". The modern world is preferable to their vision.

As time passes, more and more possibilities of manifestation, which could not have survived in more stable times, start
appearing. They give the illusion of more choices, but the reality is that they just bring more confusion. That is why
there cannot be a corporate solution but rather, as Guenon pointed out, an elite needs to form to germinate the
corporate body. These elite need to understand metaphysical principles, not to reject them ``out of principle", or
better said, from a lack of principles.

Without such knowledge, how can you make any sane decision? So that is your first task. Your earlier comment motivated
me to write an article about that; you can decide which principles make sense to you.

We are not here just to criticize, nor to put anyone in a perpetual state of anger or agitation. Knowledge of the way
things work makes one calmer and happier. It is worth the trouble for those with the intellectual capacity to
accomplish it.


\hfill

\texttt{JA on 2013-10-13 at 10:24 said: }

People forget sometimes that Germanic Austria, is quite Catholic, the seat of the empire was held by the Catholic German
Hapsburg family. I think the problem is not race it's geography – the periphery of Europe was less
affected by the Tradition, emanating outward from Rome. If we do have a Tradition as a race, it is to be found in Rome,
the eternal city that had an unbroken line of emperors beginning with Augustus through Constantine to Charlemagne and
Frederick Barbarossa.

I have a lot of issues with Catholicism as a doctrine but it is my religion, I was baptised a Catholic and refuse to
apostasise.


\hfill

\texttt{Thorsten on 2013-10-13 at 14:51 said: }

JA,

I think you are right to an extent about geography being the determining factor, as opposed to race given that Austria
and southern Germany were part of the Roman Empire before it was ``Holy". But as you touched on, I think that in the
absence of the Kaiseridee Catholicism looses much of its appeal even in northern in Germany which was never under the
pre-Christian Roman Empire.

In my particular case I was born into Methodism, my last Catholic ancestor was a great grandfather who I never knew, and
all the rest of my family has been protestant for hundreds of years (New England Puritans, Episcopals, Methodists,
German Baptist Brethren, Dutch Reformed Church). So although it's ``in the blood" Catholicism, or
any other remotely traditional religion isn't even a memory of any of my living family members.
The last time my ancestors practiced the Catholic faith in a traditional hierarchical society was probably the mid to
late 18th century southwestern Germany in the last days of the HRE. I wish I knew what to make of that.


\end{sffamily}\end{footnotesize}
